\documentclass[11pt]{article}

\usepackage[margin=2.5cm]{geometry}
\usepackage{amsmath,amssymb,amsthm}
\usepackage{bm}

\title{Exercise Sheet --Convex and Nonsmooth Optimization, Lagrangian Methods, and Finance\\[0.4em]}
\author{}
\date{}

\begin{document}
\maketitle

This sheet is organized in two parts:
\begin{itemize}
  \item \textbf{Part A -- Fundamental exercises}
  \item \textbf{Part B -- Advanced and finance-oriented exercises}
\end{itemize}

The problems are designed to practice:
convexity of quadratic and composite functions, gradient and Hessian computations,
subgradients, proximal operators, basic first-order methods (gradient / subgradient / proximal gradient),
Lagrangian and dual formulations, and applications to finance such as mean--variance portfolio optimization
and regularized portfolio selection.

\bigskip
\hrule
\bigskip

\section*{Part A -- Fundamental Exercises}

\begin{enumerate}

% -----------------------------
\item \textbf{Convex quadratic and unique minimizer}

Let
\[
f(\bm{x}) = \tfrac{1}{2}\bm{x}^\top Q \bm{x} + \bm{c}^\top \bm{x} + \alpha,
\]
where $Q \in \mathbb{R}^{n\times n}$ is symmetric, $\bm{c}\in\mathbb{R}^n$, and $\alpha\in\mathbb{R}$.

\begin{enumerate}
\item Give a necessary and sufficient condition on $Q$ for $f$ to be convex.
\item Compute the gradient $\nabla f(\bm{x})$ and the Hessian $\nabla^2 f(\bm{x})$.
\item Assume $Q$ is positive definite. Show that $f$ has a unique minimizer and give it explicitly in terms of $Q$ and $\bm{c}$.
\end{enumerate}

\bigskip
% -----------------------------
\item \textbf{Subgradients of simple nonsmooth functions}

\begin{enumerate}
\item Let $f:\mathbb{R}\to\mathbb{R}$, $f(t)=|t|$. Compute the subdifferential $\partial f(t)$ for:
\[
t>0,\qquad t<0,\qquad t=0.
\]
\item Let $g:\mathbb{R}^n\to\mathbb{R}$, $g(\bm{x})=\|\bm{x}\|_1 = \sum_{i=1}^n |x_i|$.
Describe the set $\partial g(\bm{x})$ using the subgradient of $|\cdot|$ coordinate-wise.
\end{enumerate}

\bigskip
% -----------------------------
\item \textbf{Proximal operators: projection and soft-thresholding}

\begin{enumerate}
\item Let $C=[a,b]\subset\mathbb{R}$ be a closed interval and define
\[
g(x) = \iota_C(x) =
\begin{cases}
0 & \text{if } x\in C,\\
+\infty & \text{otherwise}.
\end{cases}
\]
Show that for any $\lambda>0$,
\[
\operatorname{prox}_{\lambda g}(v)
=
\arg\min_{x\in\mathbb{R}}
\left\{ \iota_C(x) + \tfrac{1}{2\lambda} (x-v)^2 \right\}
=
P_C(v),
\]
the Euclidean projection of $v$ onto $C$.
\item Let $h(\bm{x})=\lambda\|\bm{x}\|_1$ on $\mathbb{R}^n$ with $\lambda>0$.
Derive a closed-form expression of $\operatorname{prox}_{h}(\bm{v})$ and identify it with the coordinate-wise soft-thresholding operator.
\end{enumerate}

\bigskip
% -----------------------------
\item \textbf{One step of gradient descent on a quadratic}

Consider
\[
f(\bm{x}) = \tfrac{1}{2}\bm{x}^\top Q \bm{x} - \bm{b}^\top \bm{x},
\]
where $Q$ is symmetric positive definite and $\bm{b} \in \mathbb{R}^n$.
Gradient descent with constant stepsize $\alpha>0$ is defined by
\[
\bm{x}_{k+1} = \bm{x}_k - \alpha \nabla f(\bm{x}_k).
\]

\begin{enumerate}
\item Compute $\nabla f(\bm{x})$ and give the explicit iteration formula for $\bm{x}_{k+1}$ in terms of $\bm{x}_k$, $Q$, $\bm{b}$, and $\alpha$.
\item Let $0 < \lambda_{\min} \le \lambda_{\max}$ be the smallest and largest eigenvalues of $Q$.
Give a condition on $\alpha$ ensuring convergence of gradient descent to the minimizer of $f$.
\end{enumerate}

\bigskip
\item \textbf{Projected gradient in one dimension}

Consider the constrained problem
\[
\min_{x\in[0,1]} \; f(x)= (x-2)^2.
\]

\begin{enumerate}
\item Find the unconstrained minimizer of $f$ on $\mathbb{R}$.
\item Project this minimizer onto the interval $[0,1]$ and deduce the constrained minimizer.
\item Write one step of the projected gradient iteration
\[
x_{k+1}=P_{[0,1]}(x_k-\alpha f'(x_k)),
\]
for a stepsize $\alpha>0$, and explain geometrically what the projection does.
\end{enumerate}
% -----------------------------
\item \textbf{A simple composite optimization problem}

Let $a\in\mathbb{R}$ and consider the one-dimensional problem
\[
\min_{x\in\mathbb{R}} \; F(x) = \tfrac{1}{2}(x-a)^2 + \lambda |x|,\qquad \lambda>0.
\]

\begin{enumerate}
\item Show that $F$ is convex.
\item Write down the optimality condition $0 \in \partial F(x^\star)$ and solve it explicitly to obtain $x^\star$ as a function of $a$ and $\lambda$.
\end{enumerate}

\bigskip
% -----------------------------
\item \textbf{Two-asset minimum-variance portfolio (finance)}

Consider two risky assets with (estimated) covariance matrix
\[
\Sigma = \begin{pmatrix}
\sigma_1^2 & \rho \sigma_1 \sigma_2 \\
\rho \sigma_1 \sigma_2 & \sigma_2^2
\end{pmatrix},\qquad -1<\rho<1.
\]
Let the portfolio weights be $\bm{w}=(w_1,w_2)$ subject to $w_1+w_2=1$ (no risk-free asset).

\begin{enumerate}
\item Express the portfolio variance
\[
V(\bm{w}) = \bm{w}^\top \Sigma \bm{w}
\]
as a quadratic function of the single variable $w_1$.
\item Solve the constrained problem
\[
\min_{w_1,w_2} \; V(\bm{w}) \quad \text{s.t. } w_1+w_2=1
\]
using a Lagrange multiplier and give the optimal weights $\bm{w}^\star$.
\item Discuss qualitatively the effect of the correlation $\rho$ on the optimal weights (e.g., what happens as $\rho\to -1$?).
\end{enumerate}

\bigskip

\item \textbf{Least squares regression and normal equations}

We consider a simple linear regression model. We are given data points
\[
(x_1,y_1),\dots,(x_m,y_m)\in\mathbb{R}^2,
\]
and we want to fit a straight line
\[
y \approx a x + b
\]
in the least squares sense.

\begin{enumerate}

\item \textbf{Vector–matrix formulation}

Define the parameter vector
\[
\bm{\theta} =
\begin{pmatrix}
a\\[2mm]
b
\end{pmatrix},
\qquad
\bm{y} =
\begin{pmatrix}
y_1\\
\vdots\\
y_m
\end{pmatrix},
\qquad
X =
\begin{pmatrix}
x_1 & 1\\
\vdots & \vdots\\
x_m & 1
\end{pmatrix}.
\]

\begin{enumerate}
\item Show that the vector of predictions can be written as $X\bm{\theta}$.
\item Write the least squares objective function
\[
J(\bm{\theta})=\frac12\|X\bm{\theta}-\bm{y}\|^2.
\]
\end{enumerate}

\item \textbf{Convexity}

\begin{enumerate}
\item Compute the gradient $\nabla J(\bm{\theta})$.
\item Compute the Hessian matrix $\nabla^2 J(\bm{\theta})$.
\item Show that $J$ is a convex function on $\mathbb{R}^2$.
\end{enumerate}

\item \textbf{Normal equations}

\begin{enumerate}
\item Write the first-order optimality condition
\[
\nabla J(\bm{\theta}^\star)=0.
\]
\item Show that this gives the \textbf{normal equations}
\[
X^\top X\,\bm{\theta}^\star = X^\top \bm{y}.
\]
\item Assuming $X^\top X$ is invertible, solve for $\bm{\theta}^\star$ explicitly.
\end{enumerate}

\item \textbf{Interpretation}

\begin{enumerate}
\item Explain in words what $\bm{\theta}^\star$ represents in terms of “best linear fit”.
\item In the special case where all $x_i$ are equal, discuss why $X^\top X$ fails to be invertible and how this relates to the geometry of the data.
\end{enumerate}

\end{enumerate}
\end{enumerate}
\bigskip
\hrule
\bigskip

\section*{Part B -- Advanced and Finance-Oriented Exercises}

\begin{enumerate}
\setcounter{enumi}{6} % continue numbering from Part A

% -----------------------------
\item \textbf{Subgradient optimality condition}

Let $f:\mathbb{R}^n \to \mathbb{R}$ be a proper, closed, convex function.

\begin{enumerate}
\item Show that if $0 \in \partial f(\bm{x}^\star)$, then $\bm{x}^\star$ is a global minimizer of $f$.
(Hint: use the definition of subgradient at $\bm{x}^\star$.)
\item Prove the converse: if $\bm{x}^\star$ is a global minimizer of $f$, then $0 \in \partial f(\bm{x}^\star)$.
\item Interpret this result as the nonsmooth analogue of the first-order condition $\nabla f(\bm{x}^\star)=0$ in the smooth case.
\end{enumerate}

\bigskip
% -----------------------------
\item \textbf{KKT conditions for a mean--variance portfolio (finance)}

Let $\bm{w} \in \mathbb{R}^n$ be portfolio weights, $\Sigma \in \mathbb{R}^{n\times n}$ a positive definite covariance matrix,
$\bm{\mu} \in \mathbb{R}^n$ the vector of expected returns, and $\bar{r}$ a target expected return.
Consider the constrained mean--variance problem
\[
\min_{\bm{w}} \; \tfrac{1}{2}\bm{w}^\top \Sigma \bm{w}
\quad \text{s.t. } \bm{1}^\top \bm{w} = 1,\quad \bm{\mu}^\top \bm{w} \ge \bar{r},\quad \bm{w}\ge \bm{0}.
\]

\begin{enumerate}
\item Identify the problem as a convex optimization problem (briefly justify convexity).
\item Write the Lagrangian $\mathcal{L}(\bm{w},\lambda,\nu,\bm{\gamma})$ including:
\begin{itemize}
  \item $\lambda$ for the budget constraint $\bm{1}^\top \bm{w}=1$,
  \item $\nu$ for the expected return constraint $\bm{\mu}^\top \bm{w}\ge \bar{r}$,
  \item $\bm{\gamma}$ for the non-negativity constraints $\bm{w}\ge \bm{0}$.
\end{itemize}
\item Write down the Karush--Kuhn--Tucker (KKT) conditions for optimality.
\item Explain qualitatively how these conditions reflect the trade-off between variance minimization and return/short-selling constraints.
\end{enumerate}

\bigskip
% -----------------------------
\item \textbf{Dual ascent for an equality-constrained quadratic}

Consider the problem
\[
\min_{\bm{x}\in\mathbb{R}^n} \tfrac{1}{2}\bm{x}^\top Q \bm{x}
\quad \text{s.t. } A\bm{x} = \bm{b},
\]
where $Q$ is positive definite and $A$ has full row rank.

\begin{enumerate}
\item Write the (unaugmented) Lagrangian
\[
\mathcal{L}(\bm{x},\bm{\lambda}) = \tfrac{1}{2}\bm{x}^\top Q \bm{x} + \bm{\lambda}^\top (A\bm{x}-\bm{b}).
\]
\item Compute the dual function $d(\bm{\lambda}) = \inf_{\bm{x}} \mathcal{L}(\bm{x},\bm{\lambda})$ by solving the inner minimization in closed form.
\item Show that $d(\bm{\lambda})$ is a concave quadratic function and compute $\nabla d(\bm{\lambda})$.
\item Write the dual ascent iteration for $\bm{\lambda}_{k+1}$ and comment on the meaning of the update in terms of constraint violation.
\end{enumerate}

\bigskip
% -----------------------------
\item \textbf{Augmented Lagrangian and method of multipliers}

For the same problem as above,
\[
\min_{\bm{x}} \tfrac{1}{2}\bm{x}^\top Q \bm{x} \quad \text{s.t. } A\bm{x} = \bm{b},
\]
consider the augmented Lagrangian
\[
\mathcal{L}_\rho(\bm{x},\bm{\lambda})
=
\tfrac{1}{2}\bm{x}^\top Q \bm{x}
+ \bm{\lambda}^\top (A\bm{x}-\bm{b})
+ \tfrac{\rho}{2}\|A\bm{x}-\bm{b}\|^2, \qquad \rho>0.
\]

\begin{enumerate}
\item Write the method of multipliers iteration:
\[
\bm{x}_{k+1} = \arg\min_{\bm{x}} \mathcal{L}_\rho(\bm{x},\bm{\lambda}_k),\qquad
\bm{\lambda}_{k+1} = \bm{\lambda}_k + \rho (A\bm{x}_{k+1}-\bm{b}).
\]
\item Derive the optimality condition for $\bm{x}_{k+1}$ and show that it solves a linear system involving $Q + \rho A^\top A$.
\item Give an intuition why the augmented Lagrangian can lead to better numerical behavior than pure dual ascent.
\end{enumerate}

\bigskip
% -----------------------------
\item \textbf{ADMM for a regularized portfolio problem (finance)}

We recall that in many portfolio problems, one wants at the same time:
\begin{itemize}
\item to control risk through a quadratic term involving the covariance matrix,
\item to promote sparsity of positions (few active assets) through an $\ell_1$ penalty.
\end{itemize}

Here $\Sigma \in \mathbb{R}^{n\times n}$ is a positive definite covariance matrix and 
$\bm{w}^{\mathrm{ref}}$ is a reference (benchmark) portfolio. We consider the $\ell_1$–regularized
optimization problem
\[
\min_{\bm{w}\in\mathbb{R}^n}
\left\{
\tfrac{1}{2}(\bm{w}-\bm{w}^{\mathrm{ref}})^\top \Sigma (\bm{w}-\bm{w}^{\mathrm{ref}})
+ \lambda \|\bm{w}\|_1
\right\}
\quad \text{s.t. } \bm{1}^\top \bm{w} = 1.
\]

The quadratic term measures variance relative to a benchmark, while the $\ell_1$ term encourages sparse portfolios.

\medskip
\noindent\emph{Reminder: what is ADMM?}

\begin{itemize}
\item ADMM stands for \textbf{Alternating Direction Method of Multipliers}.
\item It is designed for problems of the form
\[
\min_{x,z} f(x)+g(z)\quad \text{s.t. } Ax+Bz=c,
\]
where $f$ and $g$ are convex but handled more easily \emph{separately}.
\item ADMM:
\begin{itemize}
\item splits the problem into two easier subproblems,
\item alternates minimization with respect to $x$ and $z$,
\item updates a Lagrange multiplier (dual variable).
\end{itemize}
\item It combines ideas from:
\begin{itemize}
\item dual ascent,
\item method of multipliers,
\item operator splitting.
\end{itemize}
\item It is particularly suited to:
\begin{itemize}
\item large-scale,
\item nonsmooth,
\item composite optimization problems.
\end{itemize}
\end{itemize}

\bigskip

\begin{enumerate}

\item \textbf{Rewriting the problem as a split optimization problem}

Introduce an auxiliary variable $\bm{z}$ so that the quadratic part and the $\ell_1$ part do not act on the same variable.

\begin{enumerate}
\item Define
\[
f(\bm{w}) = \tfrac{1}{2}(\bm{w}-\bm{w}^{\mathrm{ref}})^\top \Sigma (\bm{w}-\bm{w}^{\mathrm{ref}})
+ \iota_{\{\bm{1}^\top \bm{w}=1\}}(\bm{w}),
\]
where $\iota$ is the indicator function of the constraint set.
\item Define
\[
g(\bm{z})=\lambda\|\bm{z}\|_1.
\]
\item Show that the original problem is equivalent to
\[
\min_{\bm{w},\bm{z}} \; f(\bm{w}) + g(\bm{z})
\quad \text{s.t. } \bm{w}=\bm{z}.
\]
Explain in words why this separation is useful for numerical algorithms.
\end{enumerate}

\item \textbf{Scaled-form ADMM updates}

Recall that in scaled form, ADMM works on variables $(\bm{w},\bm{z},\bm{u})$ and alternates the following steps:

\begin{enumerate}
\item A $\bm{w}$–update:
\[
\bm{w}^{k+1}
=
\arg\min_{\bm{w}}
\left\{
f(\bm{w})+\frac{\rho}{2}\|\bm{w}-\bm{z}^k+\bm{u}^k\|^2
\right\}.
\]
\item A $\bm{z}$–update:
\[
\bm{z}^{k+1}
=
\arg\min_{\bm{z}}
\left\{
g(\bm{z})+\frac{\rho}{2}\|\bm{w}^{k+1}-\bm{z}+\bm{u}^k\|^2
\right\}.
\]
\item Dual (scaled) variable update:
\[
\bm{u}^{k+1}
=
\bm{u}^k+\bm{w}^{k+1}-\bm{z}^{k+1}.
\]
\end{enumerate}

You do \emph{not} need to fully simplify the minimizers.  
Simply write down the three ADMM steps clearly for the present portfolio problem.

(Hint: the $\bm{z}$–update corresponds to a soft–thresholding operator.)

\item \textbf{Why ADMM helps here}

Give a short discussion:

\begin{itemize}
\item the $\bm{w}$–step involves:
\begin{itemize}
\item a \textbf{quadratic} function
\item plus a \textbf{linear constraint} $\bm{1}^\top \bm{w}=1$,
\end{itemize}
\item the $\bm{z}$–step involves:
\begin{itemize}
\item only the $\ell_1$ norm
\item which has a simple proximal operator (soft-thresholding),
\end{itemize}
\item the algorithm alternates between:
\begin{itemize}
\item solving a quadratic program,
\item applying a coordinatewise soft-thresholding.
\end{itemize}
\end{itemize}

Explain why this splitting makes the algorithm:
\begin{itemize}
\item numerically efficient,
\item naturally parallelizable,
\item able to handle large $n$ portfolios.
\end{itemize}

\end{enumerate}

\bigskip

\item \textbf{Transaction costs and nonsmooth portfolio optimization (finance)}

Let $\Sigma$ be a positive definite covariance matrix and $\bm{w}^{\mathrm{old}}$ the current portfolio.
We consider rebalancing to a new portfolio $\bm{w}$ with proportional transaction costs:
\[
\min_{\bm{w}} 
\left\{
\tfrac{1}{2}\bm{w}^\top \Sigma \bm{w}
+ \lambda \|\bm{w}-\bm{w}^{\mathrm{old}}\|_1
\right\}
\quad \text{s.t. } \bm{1}^\top \bm{w}=1.
\]

\begin{enumerate}
\item Explain why the cost term $\|\bm{w}-\bm{w}^{\mathrm{old}}\|_1$ models proportional transaction costs.
\item Show that the objective function is convex.
\item Write the Lagrangian including the equality constraint.
\item Derive the optimality (KKT or subgradient) conditions for a minimizer $\bm{w}^\star$.
\item Discuss qualitatively:
\begin{itemize}
\item why the $\ell_1$ term leads to \emph{sparse trades} (many zero changes),
\item how this differs from classical mean–variance optimization without transaction costs.
\end{itemize}
\end{enumerate}


% -----------------------------
\item \textbf{Comparing gradient descent and subgradient methods}

Consider the smooth strongly convex function
\[
f(x)=\tfrac{1}{2}x^2
\]
and the nonsmooth convex function
\[
g(x)=|x|.
\]

\begin{enumerate}
\item Write the gradient descent iteration for $f$ with fixed stepsize $\alpha\in(0,2)$ and solve it explicitly for $x_k$ in terms of $x_0$.
\item Write the subgradient method for $g$ with stepsize sequence $\alpha_k = \frac{c}{\sqrt{k+1}}$ and a subgradient choice $s_k \in \partial g(x_k)$.
\item Discuss qualitatively the difference in convergence behavior between the two methods (speed, oscillations, sensitivity to stepsizes).
\item Relate your discussion to the general theory: linear convergence for gradient descent on smooth strongly convex functions vs $O(1/\sqrt{k})$ function value convergence for subgradient methods.
\end{enumerate}

\end{enumerate}


\end{document}